% Original PDF: source/普通高考/2009/2009新课标文(海南,宁夏).pdf, question 10.
% Program-flow topology read from the original figure: start -> input (x,h) ->
% x<0?; yes -> y=0, no -> x<1?; the latter branches to y=x/y=1;
% all three outputs merge at 输出 y -> x>=2?; yes -> end, no -> x=x+h -> x<0?.
\begin{tikzpicture}[
  >=Stealth,
  line width=0.55pt,
  line cap=round,
  line join=round,
  every node/.style={font=\normalsize,anchor=center,align=center,
    execute at begin node={\everypar{}\leftskip=0pt\rightskip=0pt\parindent=0pt
      \hangindent=0pt\hangafter=1\parshape=0}},
  flowbox/.style={draw,minimum height=.70cm,inner xsep=4pt,inner ysep=2pt,
    align=center},
  terminal/.style={flowbox,rounded corners=.18cm,minimum width=1.55cm,
    minimum height=.74cm},
  io/.style={flowbox,shape=trapezium,trapezium left angle=72,
    trapezium right angle=108,minimum height=.80cm},
  decision/.style={flowbox,shape=diamond,aspect=2.65,minimum width=3.20cm,
    minimum height=.95cm,inner sep=1pt},
  flowarrow/.style={->,line width=0.55pt},
]
  \node[terminal] (start) at (0,11.20) {开始};
  \node[io,minimum width=3.20cm] (input) at (0,9.75) {输入 $x,h$};
  \node[decision] (testneg) at (0,8.25) {$x<0$};
  \node[decision,minimum width=3.55cm] (testone) at (2.95,6.90) {$x<1$};
  \node[flowbox,minimum width=2.05cm] (yzero) at (-3.00,5.25) {$y=0$};
  \node[flowbox,minimum width=2.05cm] (yx) at (1.35,5.25) {$y=x$};
  \node[flowbox,minimum width=2.05cm] (yone) at (4.20,5.25) {$y=1$};
  \node[io,minimum width=2.45cm] (output) at (0,3.65) {输出 $y$};
  \node[decision,minimum width=3.45cm] (testtwo) at (0,2.00) {$x\geqslant 2$};
  \node[flowbox,minimum width=3.20cm] (increment) at (4.65,3.35) {$x=x+h$};
  \node[terminal] (finish) at (0,0.55) {结束};

  \draw[flowarrow] (start.south) -- (input.north);
  % The input trunk and the return loop share one explicit merge point.
  \coordinate (entry) at (0,9.00);
  \draw (input.south) -- (entry);
  \draw[flowarrow] (entry) -- (testneg.north);

  % x<0? branches: yes to y=0; no to the second decision.
  \draw[flowarrow] (testneg.west) -- (-3.00,8.25) -- (yzero.north);
  \node[anchor=east] at (-1.35,8.42) {是};
  \draw[flowarrow] (testneg.east) -- (2.95,8.25) -- (testone.north);
  \node[anchor=west] at (1.30,8.42) {否};

  % x<1? branches to y=x and y=1.
  \draw[flowarrow] (testone.west) -- (1.35,6.90) -- (yx.north);
  \node[anchor=east] at (2.15,7.08) {是};
  \draw[flowarrow] (testone.east) -- (4.20,6.90) -- (yone.north);
  \node[anchor=west] at (3.65,7.08) {否};

  % Three output branches merge before 输出 y.
  \draw[flowarrow] (yzero.south) -- (-3.00,4.45);
  \draw (-3.00,4.45) -- (0,4.45);
  \draw[flowarrow] (yx.south) -- (1.35,4.45);
  \draw (1.35,4.45) -- (0,4.45);
  \draw[flowarrow] (yone.south) -- (4.20,4.45);
  \draw (4.20,4.45) -- (0,4.45);
  \draw[flowarrow] (0,4.45) -- (output.north);

  \draw[flowarrow] (output.south) -- (testtwo.north);
  \draw[flowarrow] (testtwo.south) -- (finish.north);
  \node[anchor=west] at (0.18,1.28) {是};

  % No branch loops upward and re-enters the first decision.
  \draw[flowarrow] (testtwo.east) -- (6.00,2.00) -- (6.00,2.95)
    -- (increment.south);
  \node[anchor=west] at (1.85,2.28) {否};
  \draw[flowarrow] (increment.north) -- (6.00,3.75) -- (6.00,9.00)
    -- (entry);

  \path[use as bounding box] (-4.35,0.18) rectangle (6.15,11.55);
\end{tikzpicture}
